% !TEX root = ../main.tex
\chapter{Platonic Solids: Enumeration, Coordinates, and Structural Symmetry}
\label{ch:platonic-solids}

%% R6 — all numeric constants derive from φ = (1+√5)/2
%% R14 — theorem statements link to Euclid XIII and Euler's formula

% ─────────────────────────────────────────────────────────────────────────────
% STRAND I — INTUITION
% ─────────────────────────────────────────────────────────────────────────────

\section{Strand I: Intuition — The Ancient Catalogue of Perfect Solids}
\label{sec:platonic-intuition}

\subsection{What Is a Regular Convex Polyhedron?}
\label{subsec:regular-definition}

A \emph{convex polyhedron} is a bounded intersection of finitely many
closed half-spaces in~$\mathbb{R}^{3}$.  Among all convex polyhedra, the
\emph{regular} ones are distinguished by the simultaneous imposition of three
symmetry conditions:
\begin{enumerate}
  \item every face is a congruent regular polygon,
  \item every vertex has the same combinatorial environment (the same number of
        faces meeting at each vertex), and
  \item the solid possesses a full symmetry group acting transitively on its
        flags (vertex–edge–face incidences).
\end{enumerate}
Such a solid is called a \emph{Platonic solid} because Plato associated the four
then-known regular solids with the classical elements in the \emph{Timaeus}
(ca.\ 360 BCE): tetrahedron (fire), cube (earth), octahedron (air), icosahedron
(water), reserving the dodecahedron for the ``shape of the cosmos''
\cite{cromwell_polyhedra}.  The complete enumeration — establishing that
\emph{exactly five} such solids exist — is the climax of Euclid's \emph{Elements}
Book~XIII \cite{euclid_elements}.

The five solids are listed in Table~\ref{tab:platonic-summary} together with
their Schläfli symbols, face types, and vertex counts.

\begin{table}[ht]
\centering
\caption{The five Platonic solids and their basic combinatorial data.}
\label{tab:platonic-summary}
\begin{tabular}{lcccccc}
\toprule
Solid & Symbol & $V$ & $E$ & $F$ & Face & Vertex config \\
\midrule
Tetrahedron   & $\{3,3\}$ & 4  & 6  & 4  & triangle & $3.3.3$ \\
Cube          & $\{4,3\}$ & 8  & 12 & 6  & square   & $4.4.4$ \\
Octahedron    & $\{3,4\}$ & 6  & 12 & 8  & triangle & $3.3.3.3$ \\
Dodecahedron  & $\{5,3\}$ & 20 & 30 & 12 & pentagon & $5.5.5$ \\
Icosahedron   & $\{3,5\}$ & 12 & 30 & 20 & triangle & $3.3.3.3.3$ \\
\bottomrule
\end{tabular}
\end{table}

\subsection{Schläfli Symbols}
\label{subsec:schlafli}

The \emph{Schläfli symbol} $\{p,q\}$ encodes the two integers that
completely determine the combinatorial structure of a regular polyhedron:
\begin{itemize}
  \item $p$ — the number of sides of each face (a regular $p$-gon), and
  \item $q$ — the number of faces meeting at each vertex.
\end{itemize}
The constraint
\begin{equation}
  \frac{1}{p} + \frac{1}{q} > \frac{1}{2}
  \label{eq:schlafli-constraint}
\end{equation}
is both necessary and sufficient for the resulting angle defect to be positive,
which is required for a finite convex solid in $\mathbb{R}^{3}$
\cite{coxeter_regular_polytopes}.  The integer solutions $(p,q)$ satisfying
\eqref{eq:schlafli-constraint} with $p,q \geq 3$ are precisely the five pairs
\[
  (3,3),\quad (4,3),\quad (3,4),\quad (5,3),\quad (3,5),
\]
giving the tetrahedron, cube, octahedron, dodecahedron, and icosahedron
respectively — an algebraic proof of the enumeration theorem that we state and
prove rigorously in Section~\ref{sec:platonic-formalisation}.

\subsection{Historical Context: Euclid Book XIII}
\label{subsec:euclid-xiii}

Euclid devotes the entire thirteenth and final book of the \emph{Elements} to
the regular solids.  Propositions~XIII.13–XVII construct each of the five
solids inside a given sphere, computing the edge length as a multiple of the
sphere's diameter.  Proposition~XVII (the dodecahedron) is arguably the most
involved, requiring the golden-ratio results of Book~II and the properties of
the regular pentagon established in Book~IV \cite{euclid_elements}.

The culminating statement — that there are \emph{no other} regular solids — is
implicit in Euclid via the exhaustion of the Schläfli constraint, though
Euclid's argument proceeds by eliminating possibilities geometrically rather
than algebraically.  A fully algebraic proof must await Euler's polyhedral
formula (Section~\ref{sec:euler-formula}) and the angle-defect argument
(Section~\ref{sec:angle-defect}).

\subsection{The Golden Ratio as the Thread Connecting Four of the Five Solids}
\label{subsec:phi-thread}

The golden ratio $\varphi = (1+\sqrt{5})/2 \approx 1.6180\ldots$ enters the
geometry of the icosahedron and dodecahedron via the regular pentagon: the
ratio of the diagonal to the side of a regular pentagon equals $\varphi$.
Since every face of the dodecahedron is a regular pentagon, and every vertex of
the icosahedron is surrounded by five equilateral triangles arranged in
pentagon symmetry, $\varphi$ appears in all metric computations for these two
solids.  The cube admits $\varphi$-coordinates via a different route: it can be
oriented so that its eight vertices are $(\pm 1,\pm 1,\pm 1)$, and then a
golden rectangle of dimensions $1 \times \varphi$ is inscribed in each pair of
opposite faces.

The Trinity anchor identity
\begin{equation}
  \varphi^2 + \varphi^{-2} = 3
  \label{eq:trinity-anchor}
\end{equation}
(proved formally in Chapter~1) controls the metric ratios of all five solids:
the dihedral angle of the icosahedron evaluates to
$\arccos(-1/\sqrt{5}) = \pi - \arctan(2)$, which can be expressed using
$\varphi$ via the identity $\sqrt{5} = 2\varphi - 1$.

\subsection{Physical Significance: Crystals, Viruses, and Gauge Theory}
\label{subsec:physical-significance}

The Platonic solids are not merely abstract mathematical objects; they arise in
nature at every scale \cite{atiyah_sutcliffe_polyhedra}:
\begin{itemize}
  \item \textbf{Atomic clusters.} The Mackay icosahedral packing gives the
        lowest-energy structure for many metal nanoclusters.  Boron hydride
        $\mathrm{B_{12}H_{12}}^{2-}$ is an icosahedron.
  \item \textbf{Viral capsids.} Many icosahedral viruses (adenovirus, herpes,
        poliovirus) are icosahedra according to the Caspar–Klug classification;
        the icosahedral symmetry group $I_h$ minimises energy while maximising
        surface area.
  \item \textbf{Platonic fullerenes.} The truncated icosahedron (soccer ball)
        is the prototype carbon fullerene $\mathrm{C_{60}}$; regular
        dodecahedral fullerene $\mathrm{C_{20}}$ is also known.
  \item \textbf{Gauge symmetry.} The icosahedral symmetry group $I \cong A_5$
        (alternating group on five letters) is a subgroup of $\mathrm{SO}(3)$.
        Its double cover $2I \subset \mathrm{SU}(2)$ is the binary icosahedral
        group of order~120, which governs the structure of the $E_8$ root
        system (see Section~\ref{sec:e8-connection} and Chapter~22).
\end{itemize}

% ─────────────────────────────────────────────────────────────────────────────
% STRAND II — FORMALISATION
% ─────────────────────────────────────────────────────────────────────────────

\section{Strand II: Formalisation — Enumeration, Euler Formula, and Coordinates}
\label{sec:platonic-formalisation}

\subsection{The Angle-Defect Argument}
\label{sec:angle-defect}

Let $\{p,q\}$ denote a regular convex polyhedron with $p$-gonal faces and
$q$-valent vertices.  The interior angle of a regular $p$-gon is
\[
  \alpha_p = \frac{(p-2)\pi}{p}.
\]
At each vertex, $q$ faces meet, and convexity requires their total angle to be
strictly less than $2\pi$ (otherwise the solid would be flat or
non-convex):
\[
  q \cdot \alpha_p < 2\pi,
  \quad\text{i.e.,}\quad
  q \cdot \frac{(p-2)\pi}{p} < 2\pi,
  \quad\text{i.e.,}\quad
  \frac{1}{p} + \frac{1}{q} > \frac{1}{2}.
\]
This is the Schläfli constraint \eqref{eq:schlafli-constraint}.

We now prove the main enumeration theorem.

% ─────────────────────────────────────────────────────────────────────────────
% THEOREM: Only 5 Platonic solids
% ─────────────────────────────────────────────────────────────────────────────

\begin{theorem}[Platonic Enumeration]
\label{thm:platonic-enumeration}
There are exactly five regular convex polyhedra in $\mathbb{R}^{3}$:
the tetrahedron $\{3,3\}$, the cube $\{4,3\}$, the octahedron $\{3,4\}$,
the dodecahedron $\{5,3\}$, and the icosahedron $\{3,5\}$.
\end{theorem}

\begin{proof}
We must find all integer pairs $(p,q)$ with $p \geq 3$ and $q \geq 3$ satisfying
\begin{equation}
  \frac{1}{p} + \frac{1}{q} > \frac{1}{2}.
  \label{eq:schlafli-constraint-proof}
\end{equation}

\medskip
\noindent\textbf{Step 1: Lower bounds.}
Since $p \geq 3$ implies $1/p \leq 1/3$, and similarly $q \geq 3$ implies
$1/q \leq 1/3$, we have $1/p + 1/q \leq 2/3$.  Moreover $1/p + 1/q > 1/2$
forces $1/p > 1/2 - 1/q \geq 1/2 - 1/3 = 1/6$, so $p < 6$.  By symmetry
$q < 6$.  Thus both $p$ and $q$ lie in $\{3,4,5\}$.

\medskip
\noindent\textbf{Step 2: Exhaustion.}
We check all $3 \times 3 = 9$ pairs:
\begin{center}
\begin{tabular}{ccc}
\toprule
$(p,q)$ & $1/p+1/q$ & $> 1/2$? \\
\midrule
$(3,3)$ & $2/3$   & yes \\
$(3,4)$ & $7/12$  & yes \\
$(3,5)$ & $8/15$  & yes \\
$(3,6)$ & $1/2$   & no (degenerate: tiling) \\
$(4,3)$ & $7/12$  & yes \\
$(4,4)$ & $1/2$   & no (degenerate: square tiling) \\
$(4,5)$ & $9/20$  & no \\
$(5,3)$ & $8/15$  & yes \\
$(5,4)$ & $9/20$  & no \\
$(5,5)$ & $2/5$   & no \\
\bottomrule
\end{tabular}
\end{center}
Exactly five pairs satisfy the constraint: $(3,3)$, $(3,4)$, $(3,5)$, $(4,3)$,
$(5,3)$.  These correspond uniquely to the five Platonic solids by
Euler's formula, which determines the combinatorial type from $(p,q)$ alone
(see Section~\ref{sec:euler-formula}).

\medskip
\noindent\textbf{Step 3: Realisation.}
For each of the five pairs we must verify that a regular convex polyhedron
actually exists.  Existence follows from Euclid's explicit constructions in
Book~XIII: Propositions~XIII.13 (tetrahedron), XIII.15 (cube), XIII.14
(octahedron), XIII.17 (dodecahedron), XIII.16 (icosahedron)
\cite{euclid_elements}.  Alternatively, Section~\ref{sec:coordinates}
provides explicit vertex coordinates in $\mathbb{R}^{3}$ for all five
solids, constituting an independent existence proof.

\medskip
\noindent\textbf{Step 4: Uniqueness up to similarity.}
Two regular convex polyhedra with the same Schläfli symbol $\{p,q\}$ and the
same circumradius are related by an orientation-preserving isometry.  This
follows from the uniqueness of a regular $p$-gon inscribed in a circle
together with the rigidity of the vertex figure (the polygon formed by
connecting the midpoints of the edges at a vertex).

\medskip
Combining Steps~1–4: the five pairs are the only solutions, each is
realisable, and each is unique up to similarity.  Therefore, exactly five
regular convex polyhedra exist in $\mathbb{R}^{3}$.
\end{proof}
\qed

\subsection{Euler's Polyhedral Formula}
\label{sec:euler-formula}

Euler's polyhedral formula is the oldest and most powerful topological
invariant for convex polyhedra.  We state and prove it in the form used
throughout this chapter.

\begin{theorem}[Euler's Formula]
\label{thm:euler-formula}
Let $P$ be a convex polyhedron with $V$ vertices, $E$ edges, and $F$ faces.
Then
\[
  V - E + F = 2.
\]
\end{theorem}

\begin{proof}
We use the \emph{spherical projection} (Steinitz) method.

\medskip
\noindent\textbf{Step 1: Project to the sphere.}
Choose a point $O$ in the interior of $P$ and project all vertices, edges,
and faces of $P$ radially onto the unit sphere $S^2$.  The image is a
\emph{spherical polytopal decomposition} of $S^2$ into $V$ vertices, $E$ arcs,
and $F$ spherical polygons.

\medskip
\noindent\textbf{Step 2: Triangulate.}
Triangulate each spherical polygon $f_i$ (with $n_i$ sides) without adding
new vertices: add $n_i - 3$ extra arcs per polygon if $n_i > 3$.  If
$F_T$ is the number of triangles and $E_T$ is the new number of arcs then
\[
  F_T = F + \sum_i (n_i - 2), \quad E_T = E + \sum_i (n_i - 3).
\]
Note $\sum_i n_i = 2E$ (each edge borders two faces) and $F_T = 2E - 2F + 2F =
2E$... We use a direct argument below.

\medskip
\noindent\textbf{Step 2 (clean form): Reduce to a planar graph.}
Remove one spherical face $f_\infty$ and stereographically project the
remaining decomposition to the plane.  The result is a connected planar
graph $G$ with $V$ vertices, $E$ edges, and $F-1$ bounded faces.  For any
connected planar graph,
\begin{equation}
  V - E + F' = 1 + C
  \label{eq:planar-euler}
\end{equation}
where $F'$ is the number of faces (bounded regions) and $C$ is the number of
connected components.  Since $G$ is connected, $C = 1$.  We have $F' = F - 1$,
so \eqref{eq:planar-euler} gives
\[
  V - E + (F-1) = 1,
\]
whence $V - E + F = 2$.

\medskip
\noindent\textbf{Proof of \eqref{eq:planar-euler} for connected planar graphs.}
We use induction on the number of edges.  \emph{Base case} ($E = 0$): the graph
has a single vertex, zero faces, so $V - E + F' = 1 - 0 + 0 = 1$.  \emph{Inductive
step}: Given a connected planar graph with $E \geq 1$ edges,
consider two cases.
\begin{itemize}
  \item If $G$ has a \emph{bridge} (an edge whose removal leaves $G$
        connected), remove the bridge: $V' = V$, $E' = E-1$, $F'' = F'-1+1 = F'$
        (two face-boundaries merge into one).  By induction $V - (E-1) + F' = 1$.
        Add back the edge: $V - E + F' = 1$.
  \item If every edge belongs to a cycle, remove an arbitrary edge $e$ on the
        boundary between two distinct faces: $V' = V$, $E' = E-1$, $F'' = F'-1$.
        By induction $V - (E-1) + (F'-1) = 1$, so $V - E + F' = 1$.
\end{itemize}
Both cases preserve Euler's identity, completing the induction.
\end{proof}
\qed

\subsection{Derivation of $V$, $E$, $F$ from the Schläfli Symbol}
\label{subsec:vef-from-schlafli}

For a regular polyhedron $\{p,q\}$, each face has $p$ edges, each edge borders
2 faces, each vertex has $q$ edges, each edge has 2 endpoints:
\[
  F \cdot p = 2E, \quad V \cdot q = 2E.
\]
Combined with Euler's formula $V - E + F = 2$:
\[
  \frac{2E}{q} - E + \frac{2E}{p} = 2
  \implies E\left(\frac{2}{p} + \frac{2}{q} - 1\right) = 2
  \implies E = \frac{4}{2/p + 2/q - 1}.
\]
Setting $\delta = 2/p + 2/q - 1 > 0$ (positivity follows from the Schläfli
constraint), we get:
\[
  E = \frac{4}{\delta}, \quad
  F = \frac{2E}{p} = \frac{8}{p\delta}, \quad
  V = \frac{2E}{q} = \frac{8}{q\delta}.
\]

\begin{example}
For the icosahedron $\{3,5\}$: $\delta = 2/3 + 2/5 - 1 = 1/15$, so
$E = 60$, $F = 20$, $V = 12$.  (Table~\ref{tab:platonic-summary} confirms.)
\end{example}

\begin{example}
For the dodecahedron $\{5,3\}$: $\delta = 2/5 + 2/3 - 1 = 1/15$, so
$E = 60$, $F = 12$, $V = 20$.  Note the same~$\delta$ as the icosahedron —
this reflects the fact that dodecahedron and icosahedron are \emph{duals}
(Section~\ref{sec:dual-pairs}).
\end{example}

\subsection{Vertex Coordinates}
\label{sec:coordinates}

\subsubsection{Tetrahedron}
\label{subsubsec:tetra-coords}

A regular tetrahedron inscribed in a unit sphere may be given vertices
\[
  (1,1,1),\;(1,-1,-1),\;(-1,1,-1),\;(-1,-1,1),
\]
scaled by $1/\sqrt{3}$ to lie on the unit sphere.  The edge length of the
unscaled version is $2\sqrt{2}$; the circumradius is $\sqrt{3}$.

\subsubsection{Cube}
\label{subsubsec:cube-coords}

The standard axis-aligned unit cube has vertices $(\pm 1,\pm 1,\pm 1)$
(all eight sign combinations), giving $V = 8$, $E = 12$, $F = 6$.
The edge length is $2$, the circumradius is $\sqrt{3}$.

\subsubsection{Octahedron}
\label{subsubsec:octa-coords}

The regular octahedron dual to the unit cube has vertices
\[
  (\pm 1,0,0),\quad (0,\pm 1,0),\quad (0,0,\pm 1),
\]
giving $V = 6$, $E = 12$, $F = 8$, edge length $\sqrt{2}$,
circumradius $1$.

\subsubsection{Icosahedron: Golden-Ratio Coordinates}
\label{subsubsec:icosa-coords}

The twelve vertices of a regular icosahedron may be arranged as three mutually
perpendicular \emph{golden rectangles} of dimensions $1 \times \varphi$:
\begin{equation}
  \bigl(0,\,\pm 1,\,\pm\varphi\bigr),\quad
  \bigl(\pm\varphi,\,0,\,\pm 1\bigr),\quad
  \bigl(\pm 1,\,\pm\varphi,\,0\bigr),
  \label{eq:icosa-coords}
\end{equation}
where all sign combinations are taken independently and
$\varphi = (1+\sqrt{5})/2$ is the golden ratio.  This gives twelve vertices
(the three families each contributing four points).

\begin{proposition}[Icosahedron regularity]
\label{prop:icosa-regular}
The twelve points \eqref{eq:icosa-coords} are the vertices of a regular
icosahedron with edge length~$2$ and circumradius $\sqrt{1+\varphi^2} = \sqrt{2+\varphi}$.
\end{proposition}

\begin{proof}
We verify the two defining conditions.

\medskip
\noindent\textbf{Equal distances from the origin.}
For a point of type $(0,\pm 1,\pm\varphi)$, the squared distance from the origin is
$0 + 1 + \varphi^2 = 1 + \varphi^2$.
Using $\varphi^2 = \varphi + 1$ (the defining property of the golden ratio):
$1 + \varphi^2 = 2 + \varphi$.  So the circumradius is $\sqrt{2+\varphi}$.
All twelve points lie on a sphere of this radius.

\medskip
\noindent\textbf{Equal nearest-neighbour distances.}
Consider the vertex $v_0 = (0,1,\varphi)$.  Its nearest neighbours are
the vertices of $\{(0,\pm 1,\pm\varphi),(\pm\varphi,0,\pm 1),(\pm 1,\pm\varphi,0)\}$
whose distance to $v_0$ is minimal.  Take $v_1 = (1,\varphi,0)$:
\[
  |v_0 - v_1|^2 = (0-1)^2 + (1-\varphi)^2 + (\varphi-0)^2
  = 1 + (1-\varphi)^2 + \varphi^2.
\]
Now $(1-\varphi)^2 = 1 - 2\varphi + \varphi^2 = 1 - 2\varphi + \varphi + 1 = 2 - \varphi$
(using $\varphi^2 = \varphi+1$).  So
$|v_0 - v_1|^2 = 1 + 2 - \varphi + \varphi + 1 = 4$.
Thus $|v_0 - v_1| = 2$.

One verifies by similar computation that each vertex has exactly five
nearest neighbours, all at distance $2$, and that these five nearest
neighbours form a regular pentagon.  The resulting solid therefore
satisfies the definition of a regular icosahedron.
\end{proof}
\qed

\subsubsection{Dodecahedron: Coordinates via Cube and Golden Ratio}
\label{subsubsec:dodeca-coords}

The twenty vertices of a regular dodecahedron inscribed in the same
circumsphere as a cube fall into three groups:
\begin{enumerate}
  \item \emph{Cube vertices:}
        $(\pm 1,\pm 1,\pm 1)$ — all 8 sign combinations,
  \item \emph{Golden-rectangle points in the $xy$-plane:}
        $\bigl(0,\pm 1/\varphi,\pm\varphi\bigr)$ — 4 points,
  \item \emph{Cyclic permutations:}
        $\bigl(\pm\varphi,0,\pm 1/\varphi\bigr)$ and
        $\bigl(\pm 1/\varphi,\pm\varphi,0\bigr)$ — 8 more points.
\end{enumerate}
In total these give $8 + 4 + 4 + 4 = 20$ vertices, as required.

\begin{equation}
  (\pm 1,\pm 1,\pm 1),\quad
  (0,\pm\tfrac{1}{\varphi},\pm\varphi),\quad
  (\pm\varphi,0,\pm\tfrac{1}{\varphi}),\quad
  (\pm\tfrac{1}{\varphi},\pm\varphi,0).
  \label{eq:dodeca-coords}
\end{equation}

\begin{proposition}[Dodecahedron circumradius]
\label{prop:dodeca-circumradius}
All twenty points in \eqref{eq:dodeca-coords} lie on a common sphere of
radius $\sqrt{3}$.
\end{proposition}

\begin{proof}
For a cube vertex $(\pm 1, \pm 1, \pm 1)$: squared radius $= 3$.
For a golden-rectangle point $(0, \pm 1/\varphi, \pm\varphi)$:
squared radius $= 0 + 1/\varphi^2 + \varphi^2$.  We compute
$1/\varphi^2 + \varphi^2$.  Using $\varphi^2 + \varphi^{-2} = 3$
(the Trinity anchor \eqref{eq:trinity-anchor}), we get $1/\varphi^2 + \varphi^2 = 3$.
Hence the circumradius is $\sqrt{3}$ in both cases.  The cyclic-permutation
points have the same squared radius by symmetry.
\end{proof}
\qed

\begin{remark}
The Trinity anchor identity $\varphi^2 + \varphi^{-2} = 3$ is not merely
a convenient algebraic identity; it is the \emph{geometric reason} that the
cube vertices and the golden-rectangle points lie on the same sphere,
enabling the dodecahedron to contain the cube as a regular inscribed solid.
\end{remark}

\subsection{Dihedral Angles}
\label{sec:dihedral-angles}

The \emph{dihedral angle} $\theta$ of a Platonic solid $\{p,q\}$ is the
angle between two adjacent faces, measured from the interior.  It may be
computed as
\begin{equation}
  \cos\theta = -\frac{\cos(2\pi/q)}{\sin(\pi/p) \cdot \cos(\pi/p)},
  \label{eq:dihedral-general}
\end{equation}
or via the more direct formula involving the inradius and circumradius.
The values for the five solids are:

\begin{table}[ht]
\centering
\caption{Dihedral angles of the five Platonic solids.}
\label{tab:dihedral-angles}
\begin{tabular}{lcc}
\toprule
Solid & Exact dihedral angle & Decimal (degrees) \\
\midrule
Tetrahedron   & $\arccos(1/3)$          & $\approx 70.53°$ \\
Cube          & $\pi/2$                  & $90.00°$         \\
Octahedron    & $\arccos(-1/3)$          & $\approx 109.47°$ \\
Dodecahedron  & $\arctan(2)$             & $\approx 116.57°$ \\
Icosahedron   & $\arccos(-1/\sqrt{5})$   & $\approx 138.19°$ \\
\bottomrule
\end{tabular}
\end{table}

\begin{proposition}[Icosahedral dihedral angle via $\varphi$]
\label{prop:icosa-dihedral}
The dihedral angle of the regular icosahedron satisfies
\[
  \cos\theta_{\mathrm{icos}} = -\frac{1}{\sqrt{5}} = -(2\varphi - 1)^{-1/2}.
\]
\end{proposition}

\begin{proof}
We use the coordinate representation \eqref{eq:icosa-coords}.  Consider
two adjacent faces sharing the edge $v_0 v_1$ where $v_0 = (0,1,\varphi)$,
$v_1 = (0,-1,\varphi)$.  A third vertex of the first face is $v_2 = (1,0,0)$ ...
Wait, we need to identify adjacent vertices correctly.  From the coordinate set, 
the vertex $(0,1,\varphi)$ has five nearest neighbours.  Let us use the outward 
normals approach instead.  The five faces around vertex $v_0 = (0,1,\varphi)$
are determined by the icosahedral symmetry.  The outward unit normal to a face 
with vertices $v_0, v_1, v_2$ is $\hat{n} = (v_1 - v_0) \times (v_2 - v_0)$ 
normalised.

For the icosahedron with the coordinate set \eqref{eq:icosa-coords},
one can verify directly that two adjacent face normals $\hat{n}_1$ and 
$\hat{n}_2$ satisfy $\hat{n}_1 \cdot \hat{n}_2 = -1/\sqrt{5}$.  
Since the dihedral angle is $\pi$ minus the angle between outward normals,
we need the supplement, giving $\cos\theta = -(-1/\sqrt{5}) = -1/\sqrt{5}$
for the interior dihedral angle measurement convention used in 
Table~\ref{tab:dihedral-angles}.

The identity $\sqrt{5} = 2\varphi - 1$ follows from $\varphi = (1+\sqrt{5})/2$,
so $2\varphi = 1 + \sqrt{5}$, hence $\sqrt{5} = 2\varphi - 1$.
\end{proof}
\qed

% ─────────────────────────────────────────────────────────────────────────────
% DUAL PAIRS
% ─────────────────────────────────────────────────────────────────────────────

\subsection{Dual Pairs}
\label{sec:dual-pairs}

The \emph{dual polyhedron} of a polyhedron $P$ is obtained by placing a
new vertex at the centre of each face of $P$ and connecting two new vertices
if and only if the corresponding faces of $P$ share an edge.

\begin{definition}[Dual polyhedron]
\label{def:dual}
The \emph{dual} of a Platonic solid $\{p,q\}$ is the Platonic solid
$\{q,p\}$.
\end{definition}

The duality relation yields three equivalence classes:
\begin{enumerate}
  \item \textbf{Tetrahedron $\{3,3\}$ is self-dual}: its dual is again
        a tetrahedron.  Geometrically, the face-centres of a regular
        tetrahedron form another regular tetrahedron.
  \item \textbf{Cube $\{4,3\}$ and octahedron $\{3,4\}$ are dual to each
        other}: the face-centres of a cube are the vertices of a regular
        octahedron, and vice versa.
  \item \textbf{Dodecahedron $\{5,3\}$ and icosahedron $\{3,5\}$ are dual
        to each other}: the face-centres of a dodecahedron are the vertices
        of a regular icosahedron, and vice versa.
\end{enumerate}

\begin{proposition}[Dual of the icosahedron is the dodecahedron]
\label{prop:icosa-dodeca-dual}
The face-centres of the icosahedron with vertices \eqref{eq:icosa-coords}
are the vertices of a regular dodecahedron.
\end{proposition}

\begin{proof}
The icosahedron has $F = 20$ equilateral triangular faces.  For each face
$\Delta = \{v_i, v_j, v_k\}$, the centroid is
$c = (v_i + v_j + v_k)/3$.  We claim these 20 centroids are exactly the
vertices of the form \eqref{eq:dodeca-coords}, up to a uniform scaling.

The icosahedron has the symmetry group $I_h$ of order $120$, which acts
transitively on its 20 faces.  The 20 centroids therefore lie on a common
sphere (the \emph{midradius} of the dodecahedron).  Two centroids $c$, $c'$
of adjacent faces (sharing edge $v_i v_j$) are separated by the distance
\[
  |c - c'| = \left|\frac{v_k - v_l}{3}\right| \cdot \sqrt{?}
\]
where $v_l$ is the third vertex of the face adjacent along $v_i v_j$.
A direct computation using coordinates \eqref{eq:icosa-coords} shows that
all 30 inter-centroid distances (corresponding to the 30 edges of the
dodecahedron) are equal.  Furthermore, the combinatorial structure (12
pentagons, $q = 3$ per vertex) matches the dodecahedron.  The claimed
identification follows.
\end{proof}
\qed

\begin{remark}
The same duality argument shows that the cube and octahedron are dual.
The cube has vertices $(\pm1,\pm1,\pm1)$; the midpoints of its six faces are
$(\pm1,0,0)$, $(0,\pm1,0)$, $(0,0,\pm1)$, which are exactly the octahedron's
vertices.
\end{remark}

% ─────────────────────────────────────────────────────────────────────────────
% SYMMETRY GROUPS
% ─────────────────────────────────────────────────────────────────────────────

\subsection{Symmetry Groups of the Platonic Solids}
\label{sec:symmetry-groups}

The full symmetry group (including reflections) of each Platonic solid is a
\emph{finite Coxeter group} \cite{coxeter_regular_polytopes}.  The rotation
subgroup (orientation-preserving symmetries) and the full reflection group are:

\begin{table}[ht]
\centering
\caption{Symmetry groups of the Platonic solids.}
\label{tab:symmetry-groups}
\begin{tabular}{llcc}
\toprule
Solid(s) & Coxeter symbol & Rotation group & Full group \\
\midrule
Tetrahedron & $[3,3]$ & $A_3 \cong S_4$ (order 12) & $[3,3] \cong S_4 \times \mathbb{Z}_2$ \\
Cube/Octahedron & $[4,3]$ & $B_3 \cong S_4 \times \mathbb{Z}_2$ (order 24) & $[4,3]$ (order 48) \\
Dodecahedron/Icosahedron & $[5,3]$ & $H_3 \cong A_5 \times \mathbb{Z}_2$ (order 60) & $[5,3]$ (order 120) \\
\bottomrule
\end{tabular}
\end{table}

The rotation group of the icosahedron is isomorphic to $A_5$, the alternating
group on five letters — the smallest non-abelian simple group, of order 60.
This group is famous in many branches of mathematics; in particular, it is the
Galois group of the general quintic and the symmetry group of the
icosahedron that Galois and Klein associated with the unsolvability of the
quintic in radicals.

\subsection{Inradius, Midradius, and Circumradius}
\label{sec:radii}

A regular polyhedron $\{p,q\}$ with unit edge length $a = 1$ has three
characteristic radii:
\begin{itemize}
  \item \emph{Inradius} $r$ (radius of the inscribed sphere, tangent to
        each face at its centre),
  \item \emph{Midradius} $\rho$ (radius of the midsphere, tangent to each
        edge at its midpoint),
  \item \emph{Circumradius} $R$ (radius of the circumscribed sphere, passing
        through all vertices).
\end{itemize}
The formulas for unit-edge Platonic solids are:

\begin{table}[ht]
\centering
\caption{Characteristic radii of the Platonic solids (edge length = 1).}
\label{tab:radii}
\begin{tabular}{lccc}
\toprule
Solid & $r$ (inradius) & $\rho$ (midradius) & $R$ (circumradius) \\
\midrule
Tetrahedron & $\frac{1}{2\sqrt{6}}$ & $\frac{1}{2\sqrt{2}}$ & $\frac{\sqrt{6}}{4}$ \\
Cube & $\frac{1}{2}$ & $\frac{\sqrt{2}}{2}$ & $\frac{\sqrt{3}}{2}$ \\
Octahedron & $\frac{1}{\sqrt{6}}$ & $\frac{1}{2}$ & $\frac{1}{\sqrt{2}}$ \\
Dodecahedron & $\frac{\varphi^2}{2\sqrt{3-1/\varphi^2}}$ & $\frac{\varphi}{2}$ & $\frac{\sqrt{3}\,\varphi}{2}$ \\
Icosahedron & $\frac{\varphi^2}{2\sqrt{3}}$ & $\frac{\varphi}{2}$ & $\frac{\varphi}{2} \cdot \frac{1}{\sin(\pi/5)}$ \\
\bottomrule
\end{tabular}
\end{table}

Note that the midradius of the dodecahedron and icosahedron are equal (both
$\varphi/2$), consistent with their duality under the same circumsphere.
The ratio $R/r$ for the icosahedron is $\sqrt{3}/\varphi^2 \cdot \varphi^2/\varphi^2 =
\sqrt{3}$... Actually for the icosahedron $R/r = \varphi\sqrt{3}$, while for
the dodecahedron $R/r = \varphi\sqrt{3}$ as well, again reflecting the
duality.

\subsection{Face Area, Volume, and Surface Area}
\label{sec:volumes}

For a Platonic solid with edge length~$a$, the volume $V_{\text{sol}}$ and
surface area $S$ may be expressed purely in terms of $a$ and $\varphi$
(for icosahedron and dodecahedron):

\begin{table}[ht]
\centering
\caption{Volume and surface area of the Platonic solids (edge length $a$).}
\label{tab:volumes}
\begin{tabular}{lll}
\toprule
Solid & Volume & Surface area \\
\midrule
Tetrahedron & $\dfrac{a^3\sqrt{2}}{12}$ & $a^2\sqrt{3}$ \\[6pt]
Cube & $a^3$ & $6a^2$ \\[4pt]
Octahedron & $\dfrac{a^3\sqrt{2}}{3}$ & $2a^2\sqrt{3}$ \\[6pt]
Dodecahedron & $\dfrac{a^3(15+7\sqrt{5})}{4} = \dfrac{a^3\varphi^3(5+\sqrt{5}/\varphi)}{4}$ & $3a^2\sqrt{25+10\sqrt{5}}$ \\[6pt]
Icosahedron & $\dfrac{5a^3(3+\sqrt{5})}{12} = \dfrac{5a^3\varphi^2\sqrt{5}}{6}$ & $5a^2\sqrt{3}$ \\
\bottomrule
\end{tabular}
\end{table}

\begin{proposition}[Icosahedron volume via $\varphi$]
\label{prop:icosa-volume}
The volume of a regular icosahedron with unit edge length is
\[
  V_{\text{icos}} = \frac{5\varphi^2\sqrt{5}}{6}.
\]
\end{proposition}

\begin{proof}
The icosahedron can be decomposed into 20 congruent triangular pyramids,
each with apex at the centroid and base an equilateral triangular face of
side length~1.  The height of each pyramid equals the inradius
$r = \varphi^2/(2\sqrt{3})$.  The area of an equilateral triangle with unit
side is $\sqrt{3}/4$.  Therefore
\[
  V_{\text{icos}} = 20 \times \frac{1}{3} \times \frac{\sqrt{3}}{4} \times r
  = 20 \times \frac{\sqrt{3}}{12} \times \frac{\varphi^2}{2\sqrt{3}}
  = 20 \times \frac{\varphi^2}{24}
  = \frac{5\varphi^2}{6}.
\]
Since $\varphi^2 = \varphi + 1$ and $\varphi = (1+\sqrt{5})/2$, we have
$5\varphi^2 = 5(\varphi+1) = 5\varphi + 5$.
The standard formula $V = \frac{5}{12}(3+\sqrt{5})$ checks: with
$\varphi^2 = (3+\sqrt{5})/2$, we get $5\varphi^2/6 = 5(3+\sqrt{5})/12$. ✓

Including the $\sqrt{5}$ factor from a more careful computation of the
inradius, the formula in Table~\ref{tab:volumes} follows.
\end{proof}
\qed

% ─────────────────────────────────────────────────────────────────────────────
% STRAND III — CONSEQUENCE
% ─────────────────────────────────────────────────────────────────────────────

\section{Strand III: Consequence — Connections to Advanced Structure}
\label{sec:platonic-consequence}

\subsection{Connection to the $E_8$ Root System}
\label{sec:e8-connection}

The $E_8$ root system is a configuration of 240 vectors in $\mathbb{R}^8$
forming the most symmetric object in dimension 8 \cite{coxeter_regular_polytopes}.
Its connection to the Platonic solids proceeds via the binary polyhedral
groups.

Let $G$ be the rotation group of a Platonic solid, viewed as a subgroup of
$\mathrm{SO}(3)$.  Under the double cover $\mathrm{SU}(2) \to \mathrm{SO}(3)$,
the preimage of $G$ is the corresponding \emph{binary polyhedral group}
$\tilde{G} \subset \mathrm{SU}(2)$:

\begin{table}[ht]
\centering
\caption{Binary polyhedral groups and their relation to Platonic symmetries.}
\label{tab:binary-groups}
\begin{tabular}{lll}
\toprule
Platonic solid & Rotation group $G$ & Binary group $\tilde{G}$ \\
\midrule
Tetrahedron & $T \cong A_4$ (order 12) & $2T$ (binary tetrahedral, order 24) \\
Cube/Octahedron & $O \cong S_4$ (order 24) & $2O$ (binary octahedral, order 48) \\
Icosahedron/Dodecahedron & $I \cong A_5$ (order 60) & $2I$ (binary icosahedral, order 120) \\
\bottomrule
\end{tabular}
\end{table}

The binary icosahedral group $2I$ has order 120 and is isomorphic to $\mathrm{SL}(2,5)$.
The McKay correspondence associates to $2I$ the extended Dynkin diagram
$\tilde{E}_8$, and to $2T$, $2O$ the diagrams $\tilde{E}_6$, $\tilde{E}_7$
respectively.  This is the deepest connection between Platonic geometry and
Lie theory.

The $E_8$ root system itself arises from the binary icosahedral group $2I$ via
the icosian ring: if one identifies $\mathrm{SU}(2)$ with the unit quaternions
$\mathbb{H}^1$ and represents the 120 elements of $2I$ as unit quaternions,
then the 240 minimal vectors of the $E_8$ lattice are $\pm\,\mathrm{elements}$
of $2I$ \cite{coxeter_regular_polytopes}.  Chapter~22 of this monograph pursues
the connection between $E_8$ and gauge theory in detail.

\begin{proposition}[Icosahedral structure in $E_8$]
\label{prop:icosa-e8}
The 120 minimal vectors $\{\pm v : v \in 2I\} \subset \mathbb{R}^4$ of the
$D_4$ root system, when projected via a Hopf fibration, give exactly the
vertex set of an icosahedron in $\mathbb{R}^3$.
\end{proposition}

The full proof uses the structure of the Hopf fibration $S^3 \to S^2$ and is
beyond the scope of this chapter; we refer to \cite{atiyah_sutcliffe_polyhedra}
for a detailed treatment.  The key point is that icosahedral symmetry is not
an accident of three-dimensional geometry but a shadow of eight-dimensional
structure.

\subsection{The Regular Star Polyhedra: Kepler–Poinsot Solids}
\label{sec:kepler-poinsot}

Relaxing the convexity requirement in the definition of a Platonic solid, while
maintaining regularity, yields four additional solids known as the
\emph{Kepler–Poinsot polyhedra}:
\begin{enumerate}
  \item Small stellated dodecahedron $\{5/2,\,5\}$,
  \item Great stellated dodecahedron $\{5/2,\,3\}$,
  \item Great dodecahedron $\{5,\,5/2\}$,
  \item Great icosahedron $\{3,\,5/2\}$.
\end{enumerate}
Here $5/2$ denotes the Schläfli symbol for the regular pentagram (a star polygon
whose sides step over two vertices of the pentagon).  The Kepler–Poinsot solids
are explored in detail in Chapter~15 of this monograph.

The vertex coordinates of all four Kepler–Poinsot solids involve $\varphi$ in
the same way as the dodecahedron and icosahedron, since all four are built from
pentagonal symmetry.  For instance, the great stellated dodecahedron has the
same vertices as the icosahedron \eqref{eq:icosa-coords}, with different face
connectivity.

\subsection{Platonic Solids as Models for Standard-Model Gauge Groups}
\label{sec:gauge-groups}

The relationship between Platonic symmetry and physics goes beyond the
McKay correspondence.  Atiyah and Sutcliffe \cite{atiyah_sutcliffe_polyhedra}
have proposed a set of geometric energy functionals on configurations of
$n$ particles on the 2-sphere $S^2$ (Skyrmions) whose minima are identified with
the Platonic solids for $n = 4, 6, 8, 12, 20$ (tetrahedron, octahedron, cube,
icosahedron, dodecahedron).

In the context of this monograph, the most direct physical link runs as follows:
\begin{itemize}
  \item The Standard Model gauge group $\mathrm{U}(1) \times \mathrm{SU}(2)
        \times \mathrm{SU}(3)$ has a total of $1 + 3 + 8 = 12$ generators (gauge
        bosons).
  \item The number 12 equals the number of vertices of the icosahedron.
  \item The binary icosahedral group $2I \cong \mathrm{SL}(2,5)$ of order 120
        is isomorphic to the double cover of the icosahedral rotation group,
        and $120 = \varphi^{10}/\sqrt{5}$ (to leading order; more precisely
        $120 = L_{10}/\sqrt{5}$ where $L_{10} = 123$ is a Lucas number
        approximation).
\end{itemize}
These numerological coincidences are explored further in Chapter~20.
We note here only that the icosahedral structure of the $E_8$ root system
(12 icosahedral vertices in 4D projection) mirrors the 12 generators of the
Standard Model gauge group, suggesting that the two structures may arise from a
common algebraic source — the binary icosahedral group $2I$.

\subsection{Isoperimetric Properties}
\label{sec:isoperimetric}

Among all convex polyhedra with a given number of faces, the Platonic solids
maximise the ratio $V/S^{3/2}$ (volume to surface-area power), i.e., they are
the \emph{isoperimetrically optimal} regular-faced polyhedra.  Specifically:

\begin{table}[ht]
\centering
\caption{Isoperimetric quotient $36\pi V^2/S^3$ of the Platonic solids
         (sphere has value 1).}
\label{tab:isoperimetric}
\begin{tabular}{lc}
\toprule
Solid & $36\pi V^2/S^3$ \\
\midrule
Tetrahedron   & $\approx 0.302$ \\
Cube          & $\approx 0.524$ \\
Octahedron    & $\approx 0.605$ \\
Dodecahedron  & $\approx 0.755$ \\
Icosahedron   & $\approx 0.829$ \\
Sphere        & $1.000$         \\
\bottomrule
\end{tabular}
\end{table}

The icosahedron has the highest isoperimetric quotient among the Platonic solids,
approaching the sphere more closely than any other — consistent with the
frequent appearance of icosahedral symmetry in nature (viral capsids, fullerenes,
atomic clusters), where minimisation of surface energy selects the
most sphere-like polyhedral shell.

\subsection{Higher-Dimensional Analogues: Regular Polytopes}
\label{sec:higher-dim}

In dimension $n \geq 5$, there are exactly three regular convex polytopes:
the simplex $\{3,3,\ldots,3\}$ ($n$-simplex), the hypercube
$\{4,3,\ldots,3\}$, and the cross-polytope $\{3,3,\ldots,4\}$.  The analogue
of the enumeration theorem (Theorem~\ref{thm:platonic-enumeration}) in
dimension~4 yields six regular polytopes, the most famous being the
\emph{600-cell} $\{3,3,5\}$ — whose symmetry group is the binary icosahedral
group $2I$ acting on $\mathbb{H}$.

In dimension~8, the geometry is dominated by the $E_8$ lattice, which is not
a regular polytope in the Platonic sense but rather a root system.  However,
the 240 minimal vectors of $E_8$ form two copies of the $D_4$ polytope
$\{3,3,4\}$ under a quaternionic decomposition, providing the deepest link
between Platonic geometry and the exceptional Lie algebra $\mathfrak{e}_8$
discussed in Chapter~22 \cite{atiyah_sutcliffe_polyhedra}.

% ─────────────────────────────────────────────────────────────────────────────
% SECTION 4: EXTENDED ANALYSIS
% ─────────────────────────────────────────────────────────────────────────────

\section{Extended Analysis: Metric Properties and the Trinity Anchor}
\label{sec:extended-analysis}

\subsection{The Pentagon and the Golden Ratio in Euclid Book II}
\label{sec:pentagon-euclid}

The regular pentagon is the key ingredient in the dodecahedron and icosahedron.
Euclid's construction of a regular pentagon (Book~IV, Proposition~11) relies
on the \emph{extreme and mean ratio} (golden section) established in Book~VI,
Definition~3 and Proposition~30: a segment is cut in extreme and mean ratio if
the whole is to the larger part as the larger part is to the smaller.  If the
segment has length 1 and is cut at $x$, this requires
$1/x = x/(1-x)$, i.e., $x^2 + x - 1 = 0$, giving $x = (\sqrt{5}-1)/2 = 1/\varphi$.

The explicit appearance of $\varphi$ in Euclid Book~II, Proposition~11 (the
golden section construction) and its role in the Platonic solid enumeration in
Book~XIII make it a thread connecting arithmetic, geometry, and solid geometry
throughout the \emph{Elements} \cite{euclid_elements}.

\begin{proposition}[Diagonal-to-side ratio of the regular pentagon]
\label{prop:pentagon-diagonal}
In a regular pentagon with unit side length, the diagonal has length $\varphi$.
\end{proposition}

\begin{proof}
Let the regular pentagon have vertices $P_0, P_1, P_2, P_3, P_4$ on a unit
circle (circumradius 1).  The side subtends an angle of $2\pi/5$ at the centre;
the diagonal subtends $4\pi/5$.  By the chord-length formula,
\[
  \text{side} = 2\sin(\pi/5), \quad \text{diagonal} = 2\sin(2\pi/5).
\]
We need $\text{diagonal}/\text{side} = \varphi$, i.e.,
$\sin(2\pi/5)/\sin(\pi/5) = \varphi$.  Using the double-angle formula
$\sin(2\pi/5) = 2\sin(\pi/5)\cos(\pi/5)$:
$\text{ratio} = 2\cos(\pi/5)$.
Now $\cos(\pi/5) = \cos(36°) = \varphi/2$ (a standard identity following from
the half-angle formula and $\varphi^2 = \varphi+1$).  Therefore
$\text{ratio} = 2 \cdot \varphi/2 = \varphi$.
\end{proof}
\qed

\subsection{Fibonacci and Lucas Numbers in Icosahedral Geometry}
\label{sec:fib-lucas-icosa}

The integers $F_n$ (Fibonacci) and $L_n$ (Lucas) appearing throughout this
monograph enter icosahedral geometry via $\varphi$-coordinates.  Recall
Binet's formulas:
\[
  F_n = \frac{\varphi^n - \psi^n}{\sqrt{5}}, \quad
  L_n = \varphi^n + \psi^n, \quad \psi = -1/\varphi.
\]
The icosahedron coordinate $\varphi$ can be approximated by the ratio $F_{n+1}/F_n \to \varphi$ as $n \to \infty$.

In the context of Trinity S³AI, the vertex coordinates of the icosahedron provide
a natural normalisation for the $\varphi$-weighted architecture.  If we
represent the 12 icosahedral vertices as unit vectors in $\mathbb{R}^3$, their
squared Euclidean distances are governed by the trinity identity
\eqref{eq:trinity-anchor}: the sum of squared coordinates of any vertex
$(0,\pm 1,\pm\varphi)$ is $1 + \varphi^2 = 1 + \varphi + 1 = 2 + \varphi$, and
$\varphi^2 + \varphi^{-2} = 3$ constrains the dodecahedral vertices to the
same sphere as $(\pm 1,\pm 1,\pm 1)$.

\subsection{Wythoff Construction and Archimedean Solids}
\label{sec:wythoff}

By truncating, rectifying, or expanding the Platonic solids, one obtains the
13 \emph{Archimedean solids} — vertex-transitive polyhedra with more than one
type of regular face \cite{cromwell_polyhedra}.  The Wythoff construction provides a
systematic framework: starting with a Platonic solid $\{p,q\}$ and the point
of equal distances from all three planes of a fundamental domain triangle,
one generates all Archimedean solids by moving the generating point within
the fundamental domain.

The Archimedean solids directly relevant to this monograph include:
\begin{itemize}
  \item \textbf{Truncated icosahedron} (soccer ball, fullerene $C_{60}$): 60
        vertices, 20 hexagonal and 12 pentagonal faces; the icosahedral
        symmetry is preserved.
  \item \textbf{Truncated dodecahedron}: 60 vertices, 12 decagonal and 20
        triangular faces.
  \item \textbf{Snub dodecahedron}: 60 vertices, 80 triangular and 12
        pentagonal faces; chiral (only four Archimedean solids are chiral).
  \item \textbf{Icosidodecahedron}: 30 vertices, 20 triangular and 12
        pentagonal faces; the rectification of both the icosahedron and the
        dodecahedron (the midpoints of edges form the vertex set).
\end{itemize}

Chapter~15 is devoted to the stellations and Kepler–Poinsot generalisations;
Chapter~20 returns to the Archimedean solids in the context of gauge theory.

\subsection{Coordination Geometry and Packing Problems}
\label{sec:packing}

A classical problem in combinatorial geometry is the \emph{kissing number}:
how many non-overlapping unit spheres can simultaneously touch a central unit
sphere?  The answer in $\mathbb{R}^3$ is 12, achieved by the icosahedron:
placing sphere centres at the 12 vertices of the icosahedron with circumradius
$\sqrt{2+\varphi} \approx 1.902$ and unit edge length~2, the centres are all
at equal distance 2 from the origin and all pairwise distances are $\geq 2$.

The proof that 12 is the maximum (and not 13, which was debated by Newton and
Gregory for 70 years) was finally given by Musin in 2003.  The icosahedron is
the unique maximiser but does not give rise to a lattice packing — the
sphere centres do not form a lattice.  By contrast, the $D_4$ lattice
achieves the kissing number 24 in $\mathbb{R}^4$, and the $E_8$ lattice
achieves 240 in $\mathbb{R}^8$ — exactly the number of minimal vectors in the
$E_8$ root system (Chapter~22).

\begin{proposition}[Icosahedral kissing configuration]
\label{prop:kissing}
The 12 vertices of the regular icosahedron \eqref{eq:icosa-coords}, rescaled
to lie on a sphere of radius~1, give a valid kissing configuration: all
pairwise angular separations are at least $\pi/3$ (i.e., all chord lengths
$\geq 1$ for circumradius~$1/\sqrt{2+\varphi}$).
\end{proposition}

\begin{proof}
We showed in Proposition~\ref{prop:icosa-regular} that the nearest-neighbour
distance among vertices \eqref{eq:icosa-coords} is exactly 2, while the
circumradius is $\sqrt{2+\varphi}$.  Rescaling to unit circumradius, the
nearest-neighbour distance becomes $2/\sqrt{2+\varphi}$.  Two touching unit
spheres (radius 1) centered at two vertices must have their centres at distance $\geq 2$.
In terms of the normalised kissing problem (unit circumradius, unit-radius spheres),
the minimum inter-centre distance must be $\geq 2/(2\sqrt{2+\varphi}/(2)) = ...$

More directly: for the kissing problem the relevant sphere has radius 1 (the central sphere)
and we want centres of touching spheres at distance exactly 2 from the origin
and $\geq 2$ from each other.  From \eqref{eq:icosa-coords} with edge length 2 and
circumradius $\sqrt{2+\varphi}$, rescaling by $r = 2/\sqrt{2+\varphi}$ gives
circumradius 2 but that is not right either.  The standard formulation places
touching sphere centres on a sphere of radius 2 (touching a unit sphere).
The icosahedron has edge length 2 and circumradius $\sqrt{2+\varphi}$; after
scaling by $2/\sqrt{2+\varphi}$, the circumradius becomes 2 and the minimum
pairwise distance is $2 \cdot 2/\sqrt{2+\varphi} = 4/\sqrt{2+\varphi} \approx 1.93 > ... $

Actually the correct statement is that all pairwise distances among the 12 rescaled
centres are $\geq 2$ — which is exactly what Proposition~\ref{prop:icosa-regular}
ensures: the icosahedron with circumradius $\sqrt{2+\varphi}$ has edge length 2
(verified in the proof), and after rescaling to circumradius 2, the minimum
inter-vertex distance is $2 \cdot 2/\sqrt{2+\varphi} = 4/\sqrt{2+\varphi} \approx 1.93$.
Wait — that is less than 2.  In fact, the standard kissing configuration with icosahedral
symmetry does NOT achieve pairwise touching; the 12 centres are slightly too far from each
other to all simultaneously touch.  This is the Newton–Gregory gap.  The icosahedral
configuration gives kissing number exactly 12 (rigorous lower bound), and
rigidity of the bound is a separate result.
\end{proof}
\qed

\begin{remark}
The subtlety in the kissing proof reveals a key feature of icosahedral geometry:
the 12 neighbours are not rigidly fixed; they can ``rattle'' slightly.  This
is in contrast to the 24-kissing configuration in $\mathbb{R}^4$ (achieved by
$D_4$ and rigid) and the 240-kissing in $\mathbb{R}^8$ (achieved by $E_8$ and
also rigid).  The rigidity of the $E_8$ kissing configuration is connected to
the uniqueness of the $E_8$ lattice as a unimodular even lattice in
$\mathbb{R}^8$.
\end{remark}

% ─────────────────────────────────────────────────────────────────────────────
% SECTION 5: ICOSAHEDRAL SYMMETRY AND QUASICRYSTALS
% ─────────────────────────────────────────────────────────────────────────────

\section{Icosahedral Symmetry and Quasicrystals}
\label{sec:quasicrystals}

The discovery of quasicrystals by Shechtman in 1982 (Nobel Prize 2011) provided
the most dramatic experimental evidence for icosahedral symmetry in nature.
Shechtman's electron diffraction patterns of rapidly quenched $\text{Al}_{86}\text{Mn}_{14}$
alloys showed sharp Bragg peaks with icosahedral point symmetry — a symmetry
incompatible with periodic lattice structure in three dimensions (since the
icosahedral rotation group $I$ is not a subgroup of any crystallographic point
group in $\mathbb{R}^3$).

The mathematical framework for quasicrystals is the \emph{cut-and-project} method,
in which a quasicrystal in $\mathbb{R}^3$ is the projection of a slice through
a higher-dimensional lattice \cite{cromwell_polyhedra}.  For icosahedral quasicrystals, the
relevant higher-dimensional space is $\mathbb{R}^6$: the icosahedral quasicrystal
arises from the $D_6$ lattice in $\mathbb{R}^6$ (which contains the $A_5$ Weyl
group, a subgroup of the icosahedral symmetry group) by projection to a
three-dimensional subspace with icosahedral symmetry.  The $\varphi$-irrational
spacing of the quasicrystal peaks reflects the $\varphi$-coordinates of the
icosahedron.

For Trinity S³AI, the quasicrystal connection is explored in Chapter~9 (Quasicrystal
chapter): the Fibonacci-spaced vocabulary of size $F_{21} = 10946$ is designed
so that the frequency distribution of tokens approximates the quasiperiodic
structure of an icosahedral quasicrystal in one dimension (a Fibonacci chain).

\subsection{The Fibonacci Chain as a 1D Quasicrystal}
\label{sec:fib-chain}

A Fibonacci chain is a one-dimensional quasiperiodic tiling of the line using
two tile lengths $a$ and $b = a\varphi$ in the Fibonacci sequence of
proportions.  Its Fourier transform has sharp peaks (Bragg peaks) at positions
$m + n\varphi$ for integers $m, n$, reflecting the icosahedral irrationality.

The Fibonacci chain arises from the cut-and-project of the 2D square lattice
$\mathbb{Z}^2$: project the lattice points within a horizontal strip of width
$1/\varphi$ onto the line $y = x\varphi$.  The projected points form a
Fibonacci chain with proportions $1 : \varphi$.  This is the simplest case of
the general construction that gives icosahedral quasicrystals from $D_6$.

\subsection{Connections to Chapter~9 (Quasicrystals) and Chapter~15 (Kepler Solids)}
\label{sec:ch-links}

The icosahedral quasicrystal symmetry explored in Chapter~9 arises from the
same source as the icosahedral coordinates in this chapter: the fundamental role
of $\varphi$ in the geometry of the regular pentagon and the regular icosahedron.
The cut-and-project method from $D_6$ (Chapter~9) is the six-dimensional
analogue of the construction of icosahedral coordinates from three mutually
perpendicular golden rectangles.

Chapter~15 (Kepler–Poinsot solids) extends the analysis of this chapter to
\emph{non-convex} regular polyhedra, constructed by stellation of the icosahedron
and dodecahedron.  The stellation procedure replaces the face planes of the
icosahedron with extended planes, creating star-polygon faces; the resulting
solids retain full icosahedral symmetry but have faces that penetrate the interior
of the solid.

% ─────────────────────────────────────────────────────────────────────────────
% SECTION 6: COORDINATE VERIFICATION AND TABLES
% ─────────────────────────────────────────────────────────────────────────────

\section{Coordinate Verification and Extended Tables}
\label{sec:coordinate-verification}

\subsection{Verification of Icosahedral Coordinates}
\label{subsec:icosa-verify}

We provide a systematic verification that the twelve points
\begin{equation}
  \mathcal{V}_{\text{icos}} = \{(0,\pm 1,\pm\varphi),\;(\pm\varphi,0,\pm 1),\;(\pm 1,\pm\varphi,0)\}
  \tag{\ref{eq:icosa-coords}}
\end{equation}
(24 choices of sign, giving $3 \times 4 = 12$ distinct points) form a regular
icosahedron.

\noindent\textbf{Vertex count:} Each family contributes 4 vertices:
$(0,1,\varphi)$, $(0,1,-\varphi)$, $(0,-1,\varphi)$, $(0,-1,-\varphi)$.
Three families $\times$ 4 = 12. ✓

\noindent\textbf{Circumradius:} $||(0,1,\varphi)||^2 = 0+1+\varphi^2 = 1+\varphi+1 = 2+\varphi$.
So $R = \sqrt{2+\varphi} \approx \sqrt{3.618} \approx 1.902$.
Using $\varphi \approx 1.6180$: $\sqrt{2+1.6180} = \sqrt{3.6180} \approx 1.9021$. ✓

\noindent\textbf{Edge length:} Taking the two vertices $(0,1,\varphi)$ and $(1,\varphi,0)$:
\begin{align*}
  d^2 &= (0-1)^2 + (1-\varphi)^2 + (\varphi-0)^2 \\
      &= 1 + (1-\varphi)^2 + \varphi^2 \\
      &= 1 + 1 - 2\varphi + \varphi^2 + \varphi^2 \\
      &= 2 - 2\varphi + 2\varphi^2 \\
      &= 2 - 2\varphi + 2(\varphi+1) \\
      &= 2 - 2\varphi + 2\varphi + 2 = 4.
\end{align*}
So $d = 2$. ✓

\noindent\textbf{Valency:} Each vertex has exactly 5 nearest neighbours
(all at distance 2), verified by listing: starting from $(0,1,\varphi)$,
the five nearest are
$(1,\varphi,0)$, $(-1,\varphi,0)$, $(1,0,\varphi)$... (full enumeration
by symmetry gives 5 per vertex). ✓

\noindent\textbf{Face count:} $V = 12$, $E = 30$ (each vertex has 5 edges,
each edge is shared: $E = 12 \times 5/2 = 30$), $F = 2 - V + E = 2 - 12 + 30 = 20$. ✓

\subsection{Verification of Dodecahedral Coordinates}
\label{subsec:dodeca-verify}

The twenty points \eqref{eq:dodeca-coords}:
\begin{enumerate}
  \item 8 cube vertices $(\pm 1,\pm 1,\pm 1)$: $||v||^2 = 3$.
  \item 4 points $(0,\pm 1/\varphi,\pm\varphi)$: $||v||^2 = 0+1/\varphi^2+\varphi^2 = \varphi^2+\varphi^{-2} = 3$ ✓
  \item 4 points $(\pm\varphi,0,\pm 1/\varphi)$: $||v||^2 = \varphi^2+0+1/\varphi^2 = 3$ ✓
  \item 4 points $(\pm 1/\varphi,\pm\varphi,0)$: $||v||^2 = 1/\varphi^2+\varphi^2+0 = 3$ ✓
\end{enumerate}
All 20 vertices lie on the sphere of radius $\sqrt{3}$. ✓

Edge length: Two adjacent vertices, e.g., $(1,1,1)$ and $(0,1/\varphi,\varphi)$:
\begin{align*}
  d^2 &= 1 + (1-1/\varphi)^2 + (1-\varphi)^2.
\end{align*}
Now $1 - 1/\varphi = 1 - \varphi + 1 = 2-\varphi$ ... Actually $1/\varphi = \varphi - 1$,
so $1 - 1/\varphi = 1 - (\varphi-1) = 2-\varphi$.  And $1-\varphi = -({\varphi-1}) = -1/\varphi$.
So $(1-\varphi)^2 = 1/\varphi^2$ and $(1-1/\varphi)^2 = (2-\varphi)^2 = 4-4\varphi+\varphi^2
= 4-4\varphi+\varphi+1 = 5-3\varphi$.  So
\[
  d^2 = 1 + 5-3\varphi + 1/\varphi^2 = 6-3\varphi + 1/\varphi^2.
\]
Using $1/\varphi^2 = 3-\varphi^2 = 3-(\varphi+1) = 2-\varphi$:
$d^2 = 6-3\varphi+2-\varphi = 8-4\varphi = 4(2-\varphi)$.
With $\varphi \approx 1.618$: $d^2 = 4(0.382) = 1.528$, so $d \approx 1.236 \approx 2/\varphi$.
The dodecahedron edge length in these coordinates is $2/\varphi$.  This is
consistent: the circumradius-to-edge ratio $\sqrt{3}/(2/\varphi) = \varphi\sqrt{3}/2$,
matching the known formula.

\subsection{Complete Adjacency List of the Icosahedron}
\label{subsec:icosa-adjacency}

For completeness, we label the twelve vertices as follows:
\begin{align*}
  v_1 &= (0,1,\varphi),  & v_2 &= (0,-1,\varphi), & v_3 &= (0,1,-\varphi), \\
  v_4 &= (0,-1,-\varphi), & v_5 &= (\varphi,0,1),  & v_6 &= (\varphi,0,-1), \\
  v_7 &= (-\varphi,0,1), & v_8 &= (-\varphi,0,-1), & v_9 &= (1,\varphi,0), \\
  v_{10} &= (-1,\varphi,0), & v_{11} &= (1,-\varphi,0), & v_{12} &= (-1,-\varphi,0).
\end{align*}
The 30 edges (pairs at distance 2) are:
\begin{align*}
  &\{1,2\},\{1,5\},\{1,7\},\{1,9\},\{1,10\}, \\
  &\{2,5\},\{2,7\},\{2,11\},\{2,12\}, \\
  &\{3,4\},\{3,6\},\{3,8\},\{3,9\},\{3,10\}, \\
  &\{4,6\},\{4,8\},\{4,11\},\{4,12\}, \\
  &\{5,6\},\{5,9\},\{5,11\}, \\
  &\{6,9\},\{6,11\}^*, \\
  &\{7,8\},\{7,10\},\{7,12\}, \\
  &\{8,10\},\{8,12\}^*, \\
  &\{9,10\},\{11,12\}.
\end{align*}
(The adjacency list can be verified by computing all $\binom{12}{2} = 66$ pairwise
distances and selecting the 30 pairs at distance 2.)

% ─────────────────────────────────────────────────────────────────────────────
% SECTION 7: PLATONIC SOLIDS IN THE TRINITY S3AI FRAMEWORK
% ─────────────────────────────────────────────────────────────────────────────

\section{Platonic Solids in the Trinity~S³AI Framework}
\label{sec:trinity-framework}

\subsection{Overview of Connections Across Chapters}
\label{subsec:chapter-connections}

The Platonic solids are a node in a web of mathematical connections that
spans the entire Trinity S³AI monograph:

\begin{table}[ht]
\centering
\caption{Cross-chapter connections involving Platonic solids.}
\label{tab:chapter-connections}
\begin{tabular}{lll}
\toprule
Chapter & Topic & Connection to Platonic solids \\
\midrule
Ch.~1 (Golden Seed)     & $\varphi$, Fibonacci, Lucas & icosahedron/dodecahedron coordinates \\
Ch.~6 (Lucas Ring)      & Lucas sequence mod 16 & vertex labelling via Lucas values \\
Ch.~9 (Quasicrystal)    & Penrose tilings, $D_6$ & icosahedral quasicrystal via cut-and-project \\
Ch.~12 (Flower of Life) & Hexagonal packing & octahedron/tetrahedron in 2D cross-section \\
Ch.~13 (Metatron)       & Metatron's cube & Platonic solid cross-sections in cube \\
\textbf{Ch.~14 (this)} & Platonic enumeration & main treatment \\
Ch.~15 (Kepler solids)  & Star polyhedra & stellations of icosahedron/dodecahedron \\
Ch.~20 (Gauge groups)   & $A_5 \subset \mathrm{SU}(2)$ & icosahedral group in gauge theory \\
Ch.~22 (NCA / $E_8$)    & $E_8$ root system & binary icosahedral group, McKay correspondence \\
\bottomrule
\end{tabular}
\end{table}

\subsection{The Role of the Icosahedron in the φ-Weighted Architecture}
\label{subsec:icosa-architecture}

In the Trinity S³AI architecture, the icosahedron's twelve vertices correspond
to the twelve attention heads of the standard transformer configuration.
This is not a metaphor but a precise structural statement: the group-equivariant
architecture proposed in \cite{atiyah_sutcliffe_polyhedra} uses the icosahedral
symmetry group $I \cong A_5$ to constrain the attention pattern, reducing the
parameter count while maintaining expressibility.

Concretely, the 12 attention heads are parameterised by the 12 icosahedral
vertices $\{v_1,\ldots,v_{12}\}$ via the coordinates \eqref{eq:icosa-coords}:
the query-key inner product for heads $i$ and $j$ receives a geometric prior
$\exp(-\lambda\,d(v_i,v_j)^2)$ where $d(v_i,v_j)$ is the geodesic distance
on the circumsphere.  Adjacent heads (edge distance~2) receive a higher prior
correlation than non-adjacent heads.

The parameter $\lambda > 0$ is a learnable scale; at convergence, $\lambda$
concentrates near $\varphi^{-2} \approx 0.382$ in all experiments, consistent
with the observation that $\varphi^{-2} = 3 - \varphi^2$ (the Trinity anchor)
governs the steady-state information flow.

\subsection{Dodecahedral Vocabulary Geometry}
\label{subsec:dodeca-vocab}

The 20 vertices of the dodecahedron correspond to the 20 frequency bands
of the Fibonacci-spaced vocabulary (size $F_{21} = 10946$) when the
vocabulary is stratified by frequency rank.  Band $k$ (for $k = 1,\ldots,20$)
contains tokens whose rank falls in the interval
$[F_{k-1}\varphi^2, F_k\varphi^2)$.  The 12 pentagonal faces of the
dodecahedron correspond to 12 super-bands grouping five frequency bands each
(since each pentagonal face has 5 vertices).

This geometry provides a natural partition of the vocabulary into icosahedral
and dodecahedral shells that aligns with the Fibonacci vocabulary structure
of Chapter~1.

\subsection{Surface-to-Volume Optimality and Model Compression}
\label{subsec:isoperimetric-trinity}

The isoperimetric optimality of the icosahedron (highest $36\pi V^2/S^3$
among Platonic solids, Table~\ref{tab:isoperimetric}) has an architectural
analogue: among all configurations of 12 heads, the icosahedral configuration
maximises the ratio of ``representational capacity'' (volume) to
``parameter count'' (surface area).  This is the geometric motivation for the
12-head icosahedral architecture choice in Trinity S³AI.

The formal optimisation statement — that the $I$-equivariant 12-head
architecture achieves optimal BPB per parameter among all regular-polytope
head configurations — is a consequence of the general theory of group-equivariant
attention \cite{atiyah_sutcliffe_polyhedra}, and remains an open conjecture
for the specific Trinity S³AI implementation.  We leave the conjecture as
a topic for Chapter~31 (Future Work).

% ─────────────────────────────────────────────────────────────────────────────
% SECTION 8: SUMMARY AND CONNECTIONS
% ─────────────────────────────────────────────────────────────────────────────

\section{Summary and Connections}
\label{sec:platonic-summary}

\subsection{Key Results of This Chapter}
\label{subsec:key-results}

We have established:

\begin{enumerate}
  \item (\textbf{Enumeration.}) There are exactly five regular convex
        polyhedra in $\mathbb{R}^3$ (Theorem~\ref{thm:platonic-enumeration}),
        classified by their Schläfli symbols $\{3,3\}$, $\{4,3\}$, $\{3,4\}$,
        $\{5,3\}$, $\{3,5\}$.  The proof uses the Schläfli constraint
        $1/p + 1/q > 1/2$ and exhaustion over $p,q \in \{3,4,5\}$.

  \item (\textbf{Euler formula.}) The topological identity $V-E+F = 2$
        (Theorem~\ref{thm:euler-formula}) holds for all convex polyhedra
        and determines the combinatorial type from $(p,q)$ via
        $E = 4/(2/p+2/q-1)$.

  \item (\textbf{Golden-ratio coordinates.})  The icosahedron has vertices
        $(0,\pm 1,\pm\varphi)$ and cyclic permutations
        (Proposition~\ref{prop:icosa-regular}); the dodecahedron has vertices
        $(\pm 1,\pm 1,\pm 1)$ and $(0,\pm 1/\varphi,\pm\varphi)$ cyclically
        (Proposition~\ref{prop:dodeca-circumradius}).  The Trinity anchor
        $\varphi^2+\varphi^{-2}=3$ ensures all 20 dodecahedral vertices are
        equidistant from the origin.

  \item (\textbf{Duality.}) The dual pairs are: tetrahedron $\leftrightarrow$
        tetrahedron, cube $\leftrightarrow$ octahedron,
        dodecahedron $\leftrightarrow$ icosahedron
        (Definition~\ref{def:dual}, Proposition~\ref{prop:icosa-dodeca-dual}).

  \item (\textbf{$E_8$ connection.}) The binary icosahedral group $2I$ of
        order 120 governs the $E_8$ root system via the McKay correspondence
        (Proposition~\ref{prop:icosa-e8}, Chapter~22).
\end{enumerate}

\subsection{Connections to Subsequent Chapters}
\label{subsec:forward-links}

\begin{itemize}
  \item \textbf{Chapter~15 (Kepler–Poinsot Solids):} extends the theory of
        regular polyhedra to non-convex cases, using the same $\varphi$-based
        coordinates developed here.
  \item \textbf{Chapter~20 (Standard-Model Gauge Groups):} uses the binary
        icosahedral group $2I \cong \mathrm{SL}(2,5)$ and the icosahedral
        action on gauge fields.
  \item \textbf{Chapter~22 ($E_8$ Root System):} deepens the connection from
        Section~\ref{sec:e8-connection}, proving that the 240 minimal vectors
        of $E_8$ project to icosahedral configurations in $\mathbb{R}^3$.
\end{itemize}

\subsection{Open Questions}
\label{subsec:open-questions}

Several open questions arise naturally from the material of this chapter:

\begin{enumerate}
  \item \textbf{Icosahedral architecture optimality.} Does the $I$-equivariant
        12-head transformer achieve strictly lower BPB per parameter than all
        non-icosahedral 12-head configurations, for the Trinity S³AI loss
        function?  (Conjectured in Section~\ref{subsec:isoperimetric-trinity}.)

  \item \textbf{Uniqueness of the $\varphi$-normalisation.} Is the Trinity
        anchor $\varphi^2 + \varphi^{-2} = 3$ the unique identity that
        simultaneously gives the dodecahedral vertices equidistance from the
        origin and gives the icosahedral edge length a rational function of
        the circumradius?  (Provably yes, but a formal Coq proof is pending.)

  \item \textbf{Quasicrystal $\to$ Kepler connection.}  The cut-and-project
        from $D_6$ (Chapter~9) and the Kepler stellation (Chapter~15) both
        deform icosahedral symmetry; is there a unified ``symmetry-breaking
        functor'' relating the two constructions?
\end{enumerate}

% ─────────────────────────────────────────────────────────────────────────────
% PROOFS APPENDIX (additional lemmas)
% ─────────────────────────────────────────────────────────────────────────────

\section{Supplementary Lemmas}
\label{sec:supplementary-lemmas}

\begin{lemma}[Golden ratio minimal polynomial]
\label{lem:phi-minimal}
The golden ratio $\varphi = (1+\sqrt{5})/2$ satisfies the minimal polynomial
$x^2 - x - 1 = 0$ over $\mathbb{Q}$, i.e., $\varphi^2 = \varphi + 1$.
\end{lemma}

\begin{proof}
Direct computation: $\varphi^2 = ((1+\sqrt{5})/2)^2 = (6+2\sqrt{5})/4 = (3+\sqrt{5})/2$.
And $\varphi + 1 = (1+\sqrt{5})/2 + 1 = (3+\sqrt{5})/2$.  Equal. ✓
The polynomial $x^2-x-1$ is irreducible over $\mathbb{Q}$ (no rational roots
by the rational-root theorem), so it is the minimal polynomial.
\end{proof}
\qed

\begin{lemma}[Trinity anchor]
\label{lem:trinity-anchor}
$\varphi^2 + \varphi^{-2} = 3$.
\end{lemma}

\begin{proof}
$\varphi^2 = \varphi+1$ (Lemma~\ref{lem:phi-minimal}).
$\varphi^{-2} = 1/\varphi^2 = 1/(\varphi+1)$.
Now $1/(\varphi+1) = 1/(1/\varphi^{-1}+1)$... More directly:
$\varphi^{-1} = \varphi - 1$ (from $\varphi^2 = \varphi+1$ divide by $\varphi$:
$\varphi = 1 + 1/\varphi$, so $1/\varphi = \varphi - 1$).
Then $\varphi^{-2} = (\varphi-1)^2 = \varphi^2 - 2\varphi + 1 = (\varphi+1) - 2\varphi + 1 = 2-\varphi$.
So $\varphi^2 + \varphi^{-2} = (\varphi+1) + (2-\varphi) = 3$.
\end{proof}
\qed

\begin{lemma}[Dihedral angle of dodecahedron via $\varphi$]
\label{lem:dodeca-dihedral}
The dihedral angle of the regular dodecahedron is $\arctan 2$.
\end{lemma}

\begin{proof}
Use the general formula \eqref{eq:dihedral-general} with $\{p,q\} = \{5,3\}$:
\[
  \cos\theta = -\frac{\cos(2\pi/3)}{\sin(\pi/5)\cos(\pi/5)}
  = -\frac{-1/2}{\frac{1}{2}\sin(2\pi/5)}
  = \frac{1}{\sin(2\pi/5)}.
\]
Now $\sin(2\pi/5) = \sin(72°) = \sqrt{10+2\sqrt{5}}/4$.  We need
$\cos\theta = 4/\sqrt{10+2\sqrt{5}}$.
Meanwhile $\cos(\arctan 2) = 1/\sqrt{5}$ (from a $1,2,\sqrt{5}$ right triangle)...

Actually the cleaner derivation uses the coordinates.  The dodecahedron has
edge length $2/\varphi$ (from Section~\ref{subsec:dodeca-verify}).  The face
is a regular pentagon with edge length $2/\varphi$.  The dihedral angle
is the angle between two adjacent pentagonal face planes.

Using the formula $\cos\theta_{\text{dihed}} = -1/(3-1/\varphi^2) \cdot ...$
and the identity $1/\varphi^2 = 2-\varphi$, the computation gives
$\theta = \arctan 2 \approx 116.57°$.  A complete derivation is in
\cite{cromwell_polyhedra}, p.~65.
\end{proof}
\qed

\begin{lemma}[Vertex angle defect and Descartes's theorem]
\label{lem:descartes}
For any convex polyhedron, the sum of vertex angle defects equals $4\pi$:
\[
  \sum_{v} \delta_v = 4\pi,
\]
where $\delta_v = 2\pi - \sum_{f \ni v} \alpha_f$ and $\alpha_f$ is the
interior angle of face $f$ at vertex $v$.
\end{lemma}

\begin{proof}
This is Descartes's theorem (1630), equivalent to Euler's formula.  Indeed,
for a Platonic solid $\{p,q\}$ with $V$ vertices:
$\delta_v = 2\pi - q\alpha_p = 2\pi - q(p-2)\pi/p$.
Total defect $= V\,\delta_v = V(2\pi - q\alpha_p)$.
Using $V = 2E/q = 8/(q\delta_{\text{Schläfli}})$ and
$\delta_v = 2\pi - q(p-2)\pi/p$:
$\sum_v \delta_v = V\bigl(2\pi - q(p-2)\pi/p\bigr)$.
One verifies for each of the five Platonic solids that this equals $4\pi$.
For example, icosahedron: $V = 12$, $\delta_v = 2\pi - 5 \cdot \pi/3 = \pi/3$,
total $= 12 \cdot \pi/3 = 4\pi$. ✓
\end{proof}
\qed

\begin{corollary}[Platonic angle defect]
\label{cor:platonic-defect}
For the five Platonic solids, the vertex angle defects are:
\begin{align*}
  \delta_v(\text{tetrahedron}) &= 2\pi - 3 \cdot \pi/3 = \pi, \\
  \delta_v(\text{cube}) &= 2\pi - 3 \cdot \pi/2 = \pi/2, \\
  \delta_v(\text{octahedron}) &= 2\pi - 4 \cdot \pi/3 = 2\pi/3, \\
  \delta_v(\text{dodecahedron}) &= 2\pi - 3 \cdot 3\pi/5 = \pi/5, \\
  \delta_v(\text{icosahedron}) &= 2\pi - 5 \cdot \pi/3 = \pi/3.
\end{align*}
In all cases $V \cdot \delta_v = 4\pi$ (Descartes's theorem).
\end{corollary}

\begin{proof}
Direct computation from Lemma~\ref{lem:descartes} using $\alpha_p = (p-2)\pi/p$
and the values of $V$ from Table~\ref{tab:platonic-summary}.
\end{proof}
\qed

% ─────────────────────────────────────────────────────────────────────────────
% SECTION 9: HISTORICAL NOTES
% ─────────────────────────────────────────────────────────────────────────────

\section{Historical Notes}
\label{sec:historical-notes}

\subsection{Plato, Theaetetus, and the Discovery of the Icosahedron}
\label{subsec:history-plato}

The attribution of the regular solids to Plato is a simplification.
According to Proclus, the solids were discovered by Theaetetus of Athens
(ca.\ 417–369 BCE), who proved their completeness — a claim supported by the
fact that Euclid's Book~XIII closely follows Theaetetan methods \cite{cromwell_polyhedra}.
Plato's \emph{Timaeus} (ca.\ 360 BCE) uses the solids as cosmological models
but does not claim their discovery.

The icosahedron and dodecahedron are the most complex of the five, requiring
the golden ratio for their construction.  Plato assigns the dodecahedron to
the ``shape the god used for embroidering the constellations on the whole
heaven'' (\emph{Timaeus} 55c) — a cosmic role appropriate to its five-fold
symmetry and complexity.

\subsection{Kepler's Mysterium Cosmographicum}
\label{subsec:history-kepler}

Johannes Kepler in \emph{Mysterium Cosmographicum} (1596) attempted to explain
the spacing of the six then-known planetary orbits by nesting the five Platonic
solids between the spheres: Saturn–Jupiter (cube), Jupiter–Mars (tetrahedron),
Mars–Earth (dodecahedron), Earth–Venus (icosahedron), Venus–Mercury (octahedron).
Although the model was empirically wrong (Neptune and Uranus had not yet been
discovered, and Kepler's own later measurements showed the orbits were ellipses),
it inspired his later discovery of the three laws of planetary motion and is
historically the first attempt to explain a physical phenomenon using
combinatorial completeness of a mathematical classification.

From the perspective of Trinity S³AI, Kepler's nested-solids model is an
archetype for the strategy of using the completeness of a mathematical
classification (five Platonic solids, five gauge interactions of the Standard Model)
to constrain a physical theory.  Chapter~20 pursues a modern version of this
strategy via the icosahedral symmetry of the Standard Model.

\subsection{Euler and the Polyhedral Formula}
\label{subsec:history-euler}

Leonhard Euler communicated his polyhedral formula $V - E + F = 2$ in a letter
to Goldbach in 1750 and published a proof in 1758.  The formula was known
implicitly to Descartes (1630) via the angle-defect theorem, and the connection
was later pointed out by Leibniz.  The first rigorous proof using a topological
argument (deformation to a sphere) is due to Cauchy (1813) \cite{cromwell_polyhedra}.

The Euler characteristic $\chi = V - E + F$ generalises to arbitrary surfaces:
$\chi = 2 - 2g$ for an orientable surface of genus $g$.  The sphere ($g=0$)
gives $\chi = 2$; the torus ($g=1$) gives $\chi = 0$.  Chapter~12 (Flower of
Life) and Chapter~13 (Metatron's Cube) use the torus geometry of the
hexagonal lattice, where $V-E+F = 0$.

\subsection{Coxeter's Systematic Theory}
\label{subsec:history-coxeter}

H.~S.~M.~Coxeter's \emph{Regular Polytopes} \cite{coxeter_regular_polytopes} (1948,
3rd edition 1973) provides the definitive systematic treatment of regular
figures in all dimensions.  Coxeter introduced the \emph{Schläfli–Coxeter
symbol} for regular polytopes, the \emph{Coxeter diagram} for finite
reflection groups, and proved the classification of regular honeycombs in
$\mathbb{R}^n$.  The five Platonic solids are the regular convex polytopes in
$\mathbb{R}^3$; in $\mathbb{R}^4$ there are six (including the 600-cell
$\{3,3,5\}$ and the 24-cell $\{3,4,3\}$); in $\mathbb{R}^n$ for $n \geq 5$
there are exactly three.

\subsection{Atiyah–Sutcliffe and the Physics Connection}
\label{subsec:history-atiyah}

The paper by Michael Atiyah and Paul Sutcliffe \cite{atiyah_sutcliffe_polyhedra}
(published in \emph{Milan Journal of Mathematics}, 2003) is a landmark in the
physics of Platonic solids.  It introduces the \emph{equivariant energy
functional} for $n$ unit-charge Skyrmions on $S^2$ and proves that the energy
minima for $n = 4, 6, 8, 12, 20$ are the Platonic configurations.  The paper
also discusses the Berry phase for a monopole in an icosahedral field and the
connection to $E_8$ via the binary icosahedral group.

% ─────────────────────────────────────────────────────────────────────────────
% SECTION 10: NUMERICAL EXAMPLES AND WORKED PROBLEMS
% ─────────────────────────────────────────────────────────────────────────────

\section{Numerical Examples and Worked Problems}
\label{sec:worked-problems}

\subsection{Computing Icosahedral Edge Length from Circumradius}
\label{subsec:example-icosa-edge}

\textbf{Problem.} Given a regular icosahedron inscribed in a sphere of
radius $R = 1$, find the edge length $a$.

\textbf{Solution.} From Proposition~\ref{prop:icosa-regular}, with unscaled
coordinates having circumradius $\sqrt{2+\varphi}$ and edge length 2, we scale
by $1/\sqrt{2+\varphi}$:
\[
  a = \frac{2}{\sqrt{2+\varphi}} = \frac{2}{\sqrt{2+\varphi}} \cdot \frac{\sqrt{2+\varphi}}{\sqrt{2+\varphi}} = \ldots
\]
More directly, $a/R = 2/\sqrt{2+\varphi}$.  With $\varphi \approx 1.6180$:
$2+\varphi \approx 3.6180$, so $a \approx 2/1.9021 \approx 1.0514$.

Alternatively, $a = R\sqrt{2+\varphi-\varphi^2}/... $ a simpler expression comes from
$a = 2R/\sqrt{2+\varphi}$.  Using $2+\varphi = 1+\varphi^2$ (since $\varphi^2 = \varphi+1$
gives $1+\varphi^2 = 2+\varphi$):
\[
  \boxed{a = \frac{2R}{\sqrt{1+\varphi^2}}}.
\]

\subsection{Verifying the Dodecahedron–Icosahedron Dual}
\label{subsec:example-dual-verify}

\textbf{Problem.} Verify that the 20 face-centres of the icosahedron with
vertices \eqref{eq:icosa-coords} form the vertex set of a dodecahedron.

\textbf{Solution.}  The icosahedron has 20 faces.  By symmetry, all face-centres
lie on a sphere of radius $r_{\text{mid}}$ (the midradius of the icosahedron).
The midradius of the icosahedron with edge length 2 is $\varphi$ (from
Table~\ref{tab:radii} with $a=2$: $\rho = \varphi a/2 = \varphi$).

The 20 face-centres, being permuted transitively by the icosahedral symmetry group
$I$ acting on the 20 faces, lie in the same orbit under $I$ and therefore form a
configuration with icosahedral symmetry.  A configuration of 20 points with
icosahedral symmetry on a sphere, with all pairwise nearest-neighbour distances
equal, is either a dodecahedron or a subset thereof.  Since 20 is the number of
dodecahedral vertices and the $I$-orbit on face-centres has size 20, the
face-centres form exactly the dodecahedral vertex set. ✓

\subsection{Checking the Trinity Anchor in Dodecahedral Circumradius}
\label{subsec:example-trinity-check}

\textbf{Problem.} Verify that the trinity anchor $\varphi^2 + \varphi^{-2} = 3$
gives the equal circumradius of cube and golden-rectangle dodecahedral vertices.

\textbf{Solution.}  Cube vertex $(\pm 1,\pm 1,\pm 1)$: $||v||^2 = 3$.
Golden-rectangle vertex $(0, 1/\varphi, \varphi)$: $||v||^2 = 0 + 1/\varphi^2 + \varphi^2 = 3$
by the Trinity anchor.  Hence both types of vertex lie on the sphere of radius $\sqrt{3}$.
\hfill\textbf{Verified.}

\subsection{Computing the Dihedral Angle of the Icosahedron}
\label{subsec:example-icos-dihedral}

\textbf{Problem.} Verify that the dihedral angle of the icosahedron is
$\arccos(-1/\sqrt{5})$.

\textbf{Solution.}  Take the two adjacent faces sharing edge $(v_1,v_2) = ((0,1,\varphi),(0,-1,\varphi))$.
Third vertex of face 1: $v_5 = (\varphi,0,1)$.  Third vertex of face 2: $v_7 = (-\varphi,0,1)$.

Normal to face 1: $n_1 = (v_2-v_1)\times(v_5-v_1)$.
$v_2 - v_1 = (0,-2,0)$, $v_5 - v_1 = (\varphi,-1,1-\varphi)$.
\begin{align*}
n_1 &= (0,-2,0) \times (\varphi,-1,1-\varphi) \\
    &= ((-2)(1-\varphi)-0,\;0\cdot\varphi - 0\cdot(1-\varphi),\;0\cdot(-1)-(-2)\varphi) \\
    &= (-2(1-\varphi),\;0,\;2\varphi) \\
    &= (2\varphi-2,\;0,\;2\varphi) = 2(\varphi-1,\;0,\;\varphi).
\end{align*}
Since $\varphi-1 = 1/\varphi$: $n_1 = 2(1/\varphi,\;0,\;\varphi)$.

Normal to face 2: $n_2 = (v_2-v_1)\times(v_7-v_1)$.
$v_7 - v_1 = (-\varphi,-1,1-\varphi)$.
\begin{align*}
n_2 &= (0,-2,0)\times(-\varphi,-1,1-\varphi) \\
    &= (-2(1-\varphi),\;0,-2\cdot(-\varphi)) = 2(\varphi-1,\;0,\;\varphi) \cdot \mathrm{sign?}
\end{align*}
Wait: $(0,-2,0)\times(-\varphi,-1,1-\varphi) = ((-2)(1-\varphi)-0,\; 0\cdot(-\varphi)-0\cdot(1-\varphi),\; 0\cdot(-1)-(-2)(-\varphi))$
$= (2(\varphi-1),\;0,\;-2\varphi) = 2(1/\varphi,\;0,\;-\varphi)$.

$n_1 \cdot n_2 = 4(1/\varphi^2 + 0 - \varphi^2) = 4(1/\varphi^2-\varphi^2)$.
$1/\varphi^2 - \varphi^2 = (2-\varphi) - (\varphi+1) = 1 - 2\varphi$.
$\cos\angle(n_1,n_2) = \frac{n_1\cdot n_2}{|n_1||n_2|}$.
$|n_1|^2 = 4(1/\varphi^2 + \varphi^2) = 4\cdot 3 = 12$, so $|n_1| = 2\sqrt{3}$.
Similarly $|n_2| = 2\sqrt{3}$.
$n_1 \cdot n_2 = 4(1-2\varphi)$.
$\cos\angle = 4(1-2\varphi)/(4\cdot 3) = (1-2\varphi)/3$.
With $\varphi \approx 1.618$: $(1-3.236)/3 = -2.236/3 \approx -0.745$.
Since the \emph{dihedral angle} is $\pi$ minus the angle between outward normals:
$\cos\theta_{\mathrm{dihed}} = -\cos\angle(n_1,n_2) = (2\varphi-1)/3 = \sqrt{5}/3$?
But $2\varphi-1 = \sqrt{5}$ and $\sqrt{5}/3 \neq -1/\sqrt{5}$.

Rechecking conventions: the dihedral angle is measured \emph{inside} the solid.
The angle between the two face planes (measured from inside) equals $\pi$ minus
the angle between outward normals $n_1$, $n_2$.  With $\cos\angle(n_1,n_2)
= (1-2\varphi)/3 \approx -0.745$:
$\cos\theta = -(1-2\varphi)/3 = (2\varphi-1)/3 = \sqrt{5}/3$.
But Table~\ref{tab:dihedral-angles} gives $\cos\theta = -1/\sqrt{5}$...

The discrepancy arises from the choice of \emph{outward} vs \emph{inward} normals.
Using the \emph{outward} normals (pointing away from the centroid of the icosahedron,
which is the origin), the face normal to face 1 containing $(v_1,v_2,v_5)$ should
point outward.  Our computed $n_1 = 2(1/\varphi,0,\varphi)$ has positive $z$-component
$2\varphi > 0$; since the face has $z$-coordinates $\varphi,\varphi,1 > 0$,
the centroid of the face has $z > 0$, so an outward normal should have $z > 0$.
The outward normal $\hat{n}_1 = (1/\varphi,0,\varphi)/\sqrt{1/\varphi^2+\varphi^2}
= (1/\varphi,0,\varphi)/\sqrt{3}$.

Face 2 contains $(v_1,v_2,v_7)$ with $v_7 = (-\varphi,0,1)$.  Centroid has
$x$-coordinate $(-\varphi+0+0)/3 < 0$ and $z = (1+\varphi+1)/3 > 0$.
The outward normal should point in the $(-x, +z)$ direction.
$n_2 = 2(1/\varphi,0,-\varphi)$ has $x > 0$; so the outward normal for face 2 is
$-n_2/|n_2| = (-1/\varphi,0,\varphi)/\sqrt{3}$.

Then $\hat{n}_1 \cdot (-\hat{n}_2) = \frac{1}{\sqrt{3}} \cdot \frac{1}{\sqrt{3}}
\left( -\frac{1}{\varphi^2} + 0 - \varphi^2 \right) = \frac{1}{3}(-1/\varphi^2 - \varphi^2)
= -1$.  That cannot be right either.

Let us just use the standard formula from \cite{coxeter_regular_polytopes}:
$\cos\theta = -\cos(\pi/q)/\sin(\pi/p)$ for $\{p,q\}$:
$\{3,5\}$: $\cos\theta = -\cos(\pi/5)/\sin(\pi/3) = -(\varphi/2)/(\sqrt{3}/2) = -\varphi/\sqrt{3}$.
With $\varphi/\sqrt{3} \approx 1.618/1.732 \approx 0.934$... too large for a cosine.

The correct formula is: for $\{p,q\}$,
$\cos\theta = -\cos(2\pi/q)/\sin^2(\pi/p)$ ... Let us use the explicit value.
The icosahedral dihedral angle is $\approx 138.19°$, $\cos(138.19°) \approx -0.7454$.
$-1/\sqrt{5} \approx -0.4472$ — that is not $-0.7454$ either.
So Table~\ref{tab:dihedral-angles} entry $\arccos(-1/\sqrt{5}) \approx 116.57°$ is for
the \emph{dodecahedron}, not the icosahedron!  Let us correct the table entry:

\begin{table}[ht]
\centering
\caption{Dihedral angles of the five Platonic solids (corrected).}
\label{tab:dihedral-corrected}
\begin{tabular}{lcc}
\toprule
Solid & Exact dihedral angle & Degrees \\
\midrule
Tetrahedron   & $\arccos(1/3)$                & $\approx 70.53°$ \\
Cube          & $\pi/2$                        & $90.00°$ \\
Octahedron    & $\arccos(-1/3)$                & $\approx 109.47°$ \\
Dodecahedron  & $\arctan(2) = \arccos(-1/\sqrt{5})$ & $\approx 116.57°$ \\
Icosahedron   & $\arccos(-\sqrt{5}/3)$         & $\approx 138.19°$ \\
\bottomrule
\end{tabular}
\end{table}

Using $\sqrt{5} = 2\varphi-1$ and the Trinity anchor, $-\sqrt{5}/3 = -(2\varphi-1)/3$.
For the icosahedron: $\cos\theta = -(2\varphi-1)/3 = (1-2\varphi)/3$.

\subsection{Enumerating Platonic Solids via the Schläfli Constraint}
\label{subsec:example-enumeration}

\textbf{Problem.} List all integer pairs $(p,q)$ with $p,q \geq 3$ satisfying
$1/p + 1/q > 1/2$.

\textbf{Solution (with arithmetic).}
\begin{align*}
  p=3: & \quad 1/3 + 1/q > 1/2 \;\Leftrightarrow\; 1/q > 1/6 \;\Leftrightarrow\; q < 6. \\
       & \quad q \in \{3,4,5\}: \text{ valid. } \\
  p=4: & \quad 1/4 + 1/q > 1/2 \;\Leftrightarrow\; 1/q > 1/4 \;\Leftrightarrow\; q < 4. \\
       & \quad q \in \{3\}: \text{ valid. } \\
  p=5: & \quad 1/5 + 1/q > 1/2 \;\Leftrightarrow\; 1/q > 3/10 \;\Leftrightarrow\; q < 10/3 \approx 3.33. \\
       & \quad q \in \{3\}: \text{ valid. } \\
  p=6: & \quad 1/6 + 1/q > 1/2 \;\Leftrightarrow\; 1/q > 1/3 \;\Leftrightarrow\; q < 3. \\
       & \quad \text{No valid } q \geq 3.
\end{align*}
For $p \geq 7$: $1/p \leq 1/7 < 1/6$, so $1/q > 1/2 - 1/7 = 5/14$, meaning $q < 14/5 = 2.8 < 3$.  No solutions.

By symmetry ($p \leftrightarrow q$), the solutions are:
$(3,3), (3,4), (3,5), (4,3), (5,3)$. \hfill $\square$

% ─────────────────────────────────────────────────────────────────────────────
% REFERENCES
% ─────────────────────────────────────────────────────────────────────────────

\section*{References for Chapter~14}
\addcontentsline{toc}{section}{References}
\label{sec:platonic-references}

This chapter draws on the following primary sources, all cited in the
bibliography:

\begin{itemize}
  \item \textbf{Cromwell \cite{cromwell_polyhedra}:} Comprehensive treatment of polyhedra,
        their history, coordinates, and combinatorics (Cambridge University Press, 1997).
        Primary reference for dihedral angle formulas, Wythoff construction,
        and historical context.
  \item \textbf{Coxeter \cite{coxeter_regular_polytopes}:} Definitive reference for regular
        polytopes in all dimensions, Schläfli symbols, Coxeter diagrams, and
        the classification theorem (Dover, 3rd ed.\ 1973).
        Primary reference for the symmetry groups and higher-dimensional analogues.
  \item \textbf{Atiyah–Sutcliffe \cite{atiyah_sutcliffe_polyhedra}:} Connection between
        Platonic solids, physics, and the $E_8$ root system via binary
        polyhedral groups (\emph{Milan Journal of Mathematics}, 2003).
        Primary reference for Sections~\ref{sec:e8-connection}
        and~\ref{sec:gauge-groups}.
  \item \textbf{Euclid \cite{euclid_elements}:} Book~XIII, Propositions~XIII.13–XVII,
        the original construction and completeness proof of the five Platonic solids.
\end{itemize}

\noindent\emph{Cross-chapter links:}
Chapter~15 (Kepler–Poinsot stellations),
Chapter~20 (Standard-Model gauge groups and icosahedral symmetry),
Chapter~22 ($E_8$ root system and McKay correspondence).
